\section{Concrete Architecture}
\label{sec:concrete}

\sectionnote{Owner: Aayan A. Kader. Target: 2.25 pages plus redrawn diagrams. Rubric: Detailed descriptions of the high-level concrete architecture as box-and-arrow diagram; Screenshots of Understand with readable fonts, edges and arrows, possibly redrawn; Good subsystem layouts.}

\begin{requirements}
  \item ``Great'' on these criteria means the architecture is at the right level, every component and dependency is discussed in terms of what it represents and why it is there, with a link to the reference architecture document on onQ, and the diagrams are carefully laid out and easy to compare with the conceptual architecture.
  \item The breakdown lists ``+ redrawn use case/sequence diagrams'' for this section, while the sequence diagrams and the use case diagram belong to Section~\ref{sec:use-cases} (D'Artagnan). Agree with D'Artagnan who redraws which diagram. This outline assumes this section redraws the Understand dependency graph.
\end{requirements}

\subsection{Box-and-Arrow Diagram}
\label{sec:concrete-diagram}

\begin{requirements}
  \item A screenshot of Understand's architecture dependency graph for the mapping in Section~\ref{sec:mapping}, exported at high resolution and cropped, with readable fonts and visible edges and arrows.
  \item A redrawn version that keeps every box where it sits in the conceptual figure (Section~\ref{sec:revised-conceptual}), with no overlapping boxes and as few edge crossings as possible. Optionally label each edge with Understand's dependency count.
  \item A legend for each figure, and text that walks the reader through it.
  \item Stay at the subsystem level: no files, classes or functions in this diagram.
\end{requirements}

\subsection{Subsystems}
\label{sec:concrete-subsystems}

\begin{requirements}
  \item For each concrete subsystem: what it represents (its crates and folders, from Section~\ref{sec:mapping}), its role in the running system, why it is there, its conceptual counterpart, and its counterpart in the reference architecture document on onQ, cited.
  \item Name the key code elements that carry each role, for example \texttt{Mux}, \texttt{Domain} and \texttt{LocalPane} for the Multiplexer.
  \item Name any subsystem with no conceptual counterpart, such as the shared protocol crate \texttt{codec} or a group of utility crates, and say how the analysis treats it.
\end{requirements}

\subsection{Dependencies}
\label{sec:concrete-dependencies}

\begin{requirements}
  \item One table with a row per edge in the figure: from, to, what the dependency represents with one example from the code (file and function), and why it exists.
  \item Mark each row as a convergence or divergence, and point to Section~\ref{sec:reflexion} for the explanation.
\end{requirements}
